\documentclass[11pt]{article}

\usepackage{amsmath,amssymb}
\usepackage{xurl}
\usepackage[hidelinks]{hyperref}
\setlength{\emergencystretch}{2em}

% Shared commands for notes in docs/paper-gaps/.

\DeclareUnicodeCharacter{2080}{\ifmmode{}_{0}\else\textsubscript{0}\fi}
\DeclareUnicodeCharacter{2081}{\ifmmode{}_{1}\else\textsubscript{1}\fi}
\DeclareUnicodeCharacter{2082}{\ifmmode{}_{2}\else\textsubscript{2}\fi}
\DeclareUnicodeCharacter{2099}{\ifmmode{}_{n}\else\textsubscript{n}\fi}

\newcommand{\Lean}{\textsc{Lean}}
\ExplSyntaxOn
\NewDocumentCommand{\leanid}{m}{
  \ifmmode
    \textsf{\detokenize{#1}}
  \else
    \group_begin:
    \urlstyle{sf}
    \tl_set:Nn \l_tmpa_tl {#1}
    \tl_replace_all:Nnn \l_tmpa_tl {\_} {_}
    \exp_args:NV \path \l_tmpa_tl
    \group_end:
  \fi
}
\ExplSyntaxOff
\providecommand{\tr}{\operatorname{tr}}
\newcommand{\tnleanIssueBase}{https://github.com/LionSR/TNLean/issues/}
\newcommand{\tnleanPrBase}{https://github.com/LionSR/TNLean/pull/}
\newcommand{\tnleanDocBase}{https://sirui-lu.com/TNLean/docs/}
\newcommand{\ghissue}[1]{\href{\tnleanIssueBase#1}{GitHub issue \##1}}
\newcommand{\ghpr}[1]{\href{\tnleanPrBase#1}{PR \##1}}
\newcommand{\doclink}[1]{\href{\tnleanDocBase#1.html}{\path{#1}}}

% Dirac notation and the identity operator, as in blueprint/src/macros/common.tex.
% \providecommand keeps note-local definitions authoritative.
\usepackage{dsfont}
\providecommand{\ket}[1]{|#1\rangle}
\providecommand{\bra}[1]{\langle #1|}
\providecommand{\braket}[2]{\langle #1 | #2 \rangle}
\providecommand{\ketbra}[2]{|#1\rangle\!\langle #2|}
\providecommand{\Id}{\mathds{1}}
\providecommand{\one}{\mathds{1}}
\providecommand{\C}{\mathbb{C}}
\providecommand{\MN}[1]{M_{#1}(\C)}
\providecommand{\MD}{M_D(\C)}

% External referencing for paper sources.  The build helper writes sanitized
% auxiliary files under build/paper-aux/ and excludes duplicated source labels.
\usepackage{xr}

% Import a generated paper auxiliary file with a caller-supplied label prefix.
% The candidate paths support builds from docs/paper-gaps/, the repository root,
% and build/paper-aux/ itself.
\newcommand{\importpaperaux}[2]{%
  \IfFileExists{../../build/paper-aux/#2.aux}{%
    \externaldocument[#1]{../../build/paper-aux/#2}%
  }{%
    \IfFileExists{build/paper-aux/#2.aux}{%
      \externaldocument[#1]{build/paper-aux/#2}%
    }{%
      \IfFileExists{#2.aux}{%
        \externaldocument[#1]{#2}%
      }{}%
    }%
  }%
}

% General external reference macros
\newcommand{\paperref}[2]{\ref{#1-#2}}
\newcommand{\papereqref}[2]{(\ref{#1-#2})}

% CPSV16 source (1606.00608 / MPDO-22-12-17-2.tex) external document and shortcuts
\importpaperaux{cpsv16-}{1606.00608/MPDO-22-12-17-2-tnlean}
\newcommand{\cpsvref}[1]{\paperref{cpsv16}{#1}}
\newcommand{\cpsveqref}[1]{\papereqref{cpsv16}{#1}}

% Verdict marker: \gapnote{<kind>}{<status>}. Kind is the policy
% classification (clarification, local-correction, scope-restriction,
% unfaithful, false-source, open-gap); status is open, wip (elimination
% underway), resolved, or historical. Machine-read by the site index;
% prints a verdict line.
\newcommand{\gapnote}[2]{\par\noindent\textbf{Verdict:} #1 (#2).\par}


\title{The Span-Element Cases in the Quantum Wielandt Inequality}
\author{}
\date{2026-04-30}

\begin{document}
\maketitle

\section{The assertion}

This note concerns the quantum Wielandt inequality of
Sanz--P\'erez-Garc\'ia--Wolf--Cirac \cite{Sanz2010Quantum}.\footnote{For
traceability, this note corresponds to
\href{https://github.com/LionSR/TNLean/issues/1049}{issue \#1049}. The local
source passage is \path{Papers/0909.5347/main.tex:360--365},
\path{Papers/0909.5347/main.tex:572--590}, and
\path{Papers/0909.5347/main.tex:645--684}.}

Let $\mathcal E_A$ be a quantum channel on $M_D(\mathbb C)$ with Kraus
operators $A_1,\ldots,A_a$. For each $n$, let
\[
  S_n(A)=\operatorname{span}\{A_{i_1}\cdots A_{i_n}:1\leq i_1,\ldots,i_n\leq a\}.
\]
Let
\[
  i(A)=\min\{n:S_n(A)=M_D(\mathbb C)\}
\]
when the set is nonempty. The source theorem uses $d$ for the number of
linearly independent Kraus operators, equivalently $d=\dim S_1(A)$ in the
minimal Kraus-span convention. With this convention Theorem~1 of the source
states that, if $\mathcal E_A$ is primitive, then $q(\mathcal E_A)\leq i(A)$
and
\[
  i(A)\leq (D^2-d+1)D^2.
\]
It also states two sharper cases:
\[
  S_1(A)\text{ contains an invertible element}
  \quad\Longrightarrow\quad
  i(A)\leq D^2-d+1,
\]
and
\[
  S_1(A)\text{ contains a non-invertible element with a nonzero eigenvalue}
  \quad\Longrightarrow\quad
  i(A)\leq D^2.
\]
These hypotheses are hypotheses about elements of the one-step span $S_1(A)$,
not only about the displayed Kraus operators themselves.

\section{The formal statements}

The formal development defines $S_n(A)$ as the exact word span and defines the
Kraus rank as
\[
  \operatorname{kr}(A)=\dim S_1(A).
\]
This is the paper's linearly independent Kraus count. Replacing the printed
symbol $d$ by $\operatorname{kr}(A)$ is therefore a faithful convention rather
than a weakening of the numerical bound.
\footnote{The formal definitions are \texttt{MPSTensor.wordSpan} in
\path{TNLean/Wielandt/SpanGrowth/CumulativeSpan.lean:57--64} and
\texttt{MPSTensor.krausRank} in
\path{TNLean/Wielandt/Primitivity/PaperDefinitions.lean:85--91}. The local
source identifies $d$ as the number of linearly independent Kraus operators in
\path{Papers/0909.5347/main.tex:360--365}; the proof of Lemma~1 uses
$\dim T_1(A)=d$ in \path{Papers/0909.5347/main.tex:579--583}.}

The general bound is formalized with the exact word-span index:
\[
  i(A)\leq (D^2-\operatorname{kr}(A)+1)D^2.
\]
The statement assumes the paper's primitive spreading condition and
trace-preserving normalization. The proof passes through eventual full Kraus
rank and uses a blocking argument, but the conclusion is the fixed-length
statement for $S_n(A)$, not merely a cumulative-span statement.
\footnote{The paper-facing general theorem is
\texttt{MPSTensor.iIndex\_le\_general\_of\_isPrimitivePaper} in
\path{TNLean/Wielandt/PaperResults/WielandtInequality.lean:251--348}. The
formal equivalence between paper primitivity and eventual full Kraus rank under
normalization is in \path{TNLean/Wielandt/Primitivity/Equivalence.lean:169--191}.}

The standalone formal statements corresponding to the two sharper cases have a
different form. The invertible case assumes that one of the listed Kraus
operators $A_{i_0}$ is invertible, and concludes
\[
  S_{D^2-\operatorname{kr}(A)+1}(A)=M_D(\mathbb C).
\]
The non-invertible case assumes that one listed Kraus operator $A_{i_0}$ is
non-invertible and has a nonzero eigenvalue, and concludes
\[
  S_{D^2}(A)=M_D(\mathbb C).
\]
Thus the formal hypotheses are hypotheses on a displayed generator, whereas the
source hypotheses are hypotheses on the whole subspace $S_1(A)$.
\footnote{The formal invertible case is
\texttt{MPSTensor.wordSpan\_eq\_top\_of\_isPrimitivePaper\_of\_isUnit} and
\texttt{MPSTensor.iIndex\_le\_of\_isPrimitivePaper\_of\_isUnit} in
\path{TNLean/Wielandt/PaperResults/WielandtInequality.lean:182--217}. The
formal non-invertible case is
\texttt{MPSTensor.wordSpan\_eq\_top\_of\_isPrimitivePaper\_of\_noninvertible\_eigenvector}
and \texttt{MPSTensor.iIndex\_le\_sq\_of\_noninvertible\_eigenvector} in
\path{TNLean/Wielandt/PaperResults/WielandtInequality.lean:119--178}.}

\section{The point at issue}

For a fixed Kraus representation, the assertion that $S_1(A)$ contains an
invertible element is weaker than the assertion that one of the matrices
$A_i$ is invertible. The same distinction applies to a non-invertible element
of $S_1(A)$ with a nonzero eigenvalue. A linear combination of Kraus operators
can have one of these properties even when no displayed Kraus operator has it.

This distinction is mathematical, not merely syntactic. In the source theorem,
the element supplied by the hypothesis may be any matrix in $S_1(A)$. The
formal theorem can be applied directly only after that element has been made
into a named one-step generator, or after the proof is generalized so that the
one-step element itself replaces the named generator.

The invertible case is expected to generalize directly: if $X\in S_1(A)$ is
invertible, then left multiplication by $X$ maps $S_n(A)$ into $S_{n+1}(A)$ and
is injective. This is the dimension-growth mechanism used in the source proof.
The formal statement instead uses this mechanism only for a listed Kraus
operator. The non-invertible case similarly should be stated for an
element $X\in S_1(A)$ with a nonzero eigenvalue, or else accompanied by a
Kraus-representation-change lemma showing that the same channel and the same
word spans may be represented with $X$ as a displayed Kraus operator.

\section{What is not a discrepancy}

The replacement of $d$ by $\operatorname{kr}(A)=\dim S_1(A)$ is not a paper
error and not a formal weakening. It records explicitly the linearly
independent Kraus count used by the source proof.

Nor is the cumulative-span material, by itself, a theorem-level deviation.
There are intermediate formal results about
\[
  T_n(A)=S_0(A)+S_1(A)+\cdots+S_n(A),
\]
including the bound $T_{D^2}(A)=M_D(\mathbb C)$ under normality. These results
are weaker than the full quantum Wielandt theorem, but the paper-facing general
theorem proves the exact word-span bound for $i(A)$.
\footnote{The cumulative-span theorem is
\texttt{MPSTensor.cumulativeSpan\_eq\_top} in
\path{TNLean/Wielandt/SpanGrowth/NonzeroTraceProduct.lean:164--197}; the
paper-facing exact-word-span theorem is the general theorem cited above. The
blueprint makes this distinction in
\path{blueprint/src/chapter/ch07_wielandt.tex:461--490} and states the general
bound in \path{blueprint/src/chapter/ch07_wielandt.tex:492--548}.}

Finally, the lower-level use of normality is not itself a mismatch. In these
formal statements, normality is eventual full Kraus rank. Under the paper's
normalization, the formal statements prove the equivalence between the paper's
primitive condition and eventual full Kraus rank, so using normality as an intermediate
hypothesis is a legitimate intermediate argument.

\section{Conclusion}

The conclusion is a provisional difference in hypotheses for the sharper cases.
The general quantum Wielandt bound is formalized with the correct fixed-length
index and with the source's linearly independent Kraus count made explicit as
$\operatorname{kr}(A)$. The formal proof through normality and blocking is
therefore not a gap in the general bound.

The standalone sharper cases are not yet stated in the same generality as
Theorem~1 of \cite{Sanz2010Quantum}. The source assumes that the one-step span
$S_1(A)$ contains an invertible element, or a non-invertible element with a
nonzero eigenvalue. The formal statements assume that such an element is one of
the listed Kraus operators. This is a strengthening of the hypotheses for those
standalone cases. To remove the deviation, the formal development should either
prove the two sharper cases for arbitrary qualifying elements of $S_1(A)$, or
prove an explicit Kraus-representation invariance theorem that reduces the
span-element hypotheses to the displayed-generator hypotheses without changing
the channel or the relevant word-span indices.

\bibliographystyle{alpha}
\bibliography{references}

\end{document}